\documentclass[12pt]{article}
\usepackage[utf8]{inputenc}
\usepackage[T1]{fontenc}
\usepackage{hyperref}
\title{Atoro Language Resources for K-12 Curriculum Development}
\author{Compiled by Claude Sonnet 5 (claude-sonnet-5) for Andr\'as Kornai}
\date{}
\begin{document}
\maketitle

\section{Introduction}
Atoro (ISO 639-3 tdv; Glottocode toro1249) is a Plateau language of Nigeria, also called Toro, Turkwam, Turkwan, Turkpa or Turkpan. Search results place the speakers in Nasarawa State, in Turkwan Mission and Kanja. Wikipedia cites about 3,900 speakers (2000), while Global Recordings Network lists a population of 11,000. Little published material was found, and nothing designed for schools.

\section{Primary Educational Resources}
No primers, textbooks or school programmes for Atoro were found in the searches.

\section{Linguistic and Reference Materials}
\begin{itemize}
\item Wikipedia, Toro language (overview, classification, speaker figures): \url{https://en.wikipedia.org/wiki/Toro_language}
\item A search summary mentions an Atoro dictionary at Turkwan Mission, but I did not see a page for it, so I give no link and cannot confirm its details.
\end{itemize}

\section{Religious and Community Resources}
Joshua Project reports no Scripture in Atoro; its Turkwam people-group page is \url{https://joshuaproject.net/people_groups/15656}. Its language page is \url{https://joshuaproject.net/index.php/languages/tdv}.

\section{Digital and Audio Resources}
Global Recordings Network lists Atoro audio (Bible stories): \url{https://globalrecordings.net/en/language/tdv}

\section{Curriculum Development Context}
The language is vital but has no known written educational materials. A curriculum would have to start from orthography development and community-led material creation. The existence and contents of the dictionary should be checked directly.
\end{document}